\section{Conclusion}
\label{sec:conclusion}

This paper presented DSPHN, a novel semantic segmentation architecture
whose primary contribution is \textbf{architectural innovation for
hazard-aware scene understanding}. Three architectural components were
proposed and validated:

\begin{enumerate}
\item \textbf{Dual-Stream Appearance-Context Decoder:} Parallel standard
and dilated convolutional decoder pathways that separately model local
visual hazard cues and spatial contextual relationships, addressing the
complementary nature of hazard recognition signals.

\item \textbf{Severity-Guided Adaptive Fusion Gates:} Learnable gates
initialised from domain injury severity scores that blend appearance
and context features, providing domain-informed fusion prior to
hazard classification.

\item \textbf{Two-Stage Hazard Refinement Mechanism:} A hierarchical
conditional design where Stage~1 produces a binary safe/unsafe map
from dual-stream fused features, and Stage~2 performs fine-grained
four-class classification conditioned on Stage~1 output as a direct
structural input.
\end{enumerate}

A controlled ablation study under identical training conditions
confirms that each component contributes measurable improvement:
DSPHN (mIoU~0.4198) outperforms the dual-stream variant alone
(0.3894, $+$0.030) and the two-stage variant alone (0.3793,
$+$0.041). This ablation provides primary evidence of the
architectural contributions independent of dataset limitations.

Severe class imbalance affected performance on rare hazard classes.
Sharp~Objects achieved zero IoU due to pixel-level sparsity of
0.06\% of test pixels — a known challenge in long-tail segmentation
that represents a dataset limitation rather than an architectural
failure, as confirmed by consistent ablation improvements across
all evaluated classes.

\noindent\textbf{Limitations.}
Three limitations are acknowledged: (1)~zero IoU on Sharp~Objects
due to extreme pixel sparsity despite class weighting and COCO
augmentation; (2)~15-epoch training on a T4~GPU limits convergence;
(3)~256$\times$256 resolution may miss fine-grained boundaries of
small hazard instances.

\noindent\textbf{Future Research Directions.}

\begin{itemize}
\item \textit{Larger hazard datasets:} A dedicated indoor hazard
benchmark with balanced representation of all hazard classes,
particularly sharp objects and wet floor surfaces, would
substantially improve rare class performance and enable more
rigorous evaluation.

\item \textit{Better rare-class balancing:} Patch-level oversampling
of hazard regions, combined with synthetic data augmentation for
rare classes, provides a more principled approach than image-level
oversampling which was found to reduce overall performance.

\item \textit{Real-time deployment:} A lightweight DSPHN variant
with a MobileNet or EfficientNet-Lite encoder targeting real-time
inference would enable deployment on edge devices in domestic and
healthcare environments.

\item \textit{Robotics and autonomous navigation:} Integration of
DSPHN into robot navigation pipelines for hazard-aware path
planning represents a direct application of this work, where the
Stage~1 binary map provides a fast hazard alert and Stage~2
provides detailed classification for navigation decisions.
\end{itemize}

DSPHN demonstrates that explicit hazard-aware architectural design
— through dual-stream feature disentanglement, severity-guided
fusion, and hierarchical safety-aware inference — can improve
indoor safety-oriented scene understanding and provides a
foundation for future research in safety-critical semantic
segmentation.

\section*{Acknowledgements}
ADE20K annotations: Zhou~\etal~\cite{zhou2017scene}.
COCO annotations: Lin~\etal~\cite{lin2014coco}.
Manual hazard annotations created using Make~Sense~AI.
Training on Google Colab NVIDIA T4 GPU.
